% MG05a. Tras el párrafo de extensión y procedencia que sigue a
% act21:tpk:logaritmo; antes de «Orden de lectura y movimiento».
\subsubsection{Calibres del observable multiplicativo y normalización del generador}
\label{mg05:calibres}

La elección del exponente discreto admite una comparación finita en la
que permanecen fijos el soporte APP, las condiciones iniciales de los
dos cursores, el selector trítico, el calendario de inversión y la regla
emisora. La variación se concentra en el observable de la hoja
multiplicativa. Sean \(\ell_{2}\) y \(\ell_{5}\) los logaritmos del
grupo \((\mathbb Z/9\mathbb Z)^\times\) con generadores \(2\) y
\(5\), respectivamente, extendidos por cero sobre \(\{3,6,9\}\).
Se comparan con \(v_3\), evaluada sobre los representantes positivos
\(1,\ldots,9\), y con la lectura identidad sobre esos mismos
representantes.

\begin{proposition}[Normalización del logaritmo multiplicativo]
\label{mg05:normalizacion}
Los dos isomorfismos del grupo de unidades módulo nueve con
\(\mathbb Z/6\mathbb Z\), normalizados por enviar un generador a
\(1\), son \(\ell_2\) y \(\ell_5\). Satisfacen
\[
 \ell_5(x)\equiv-\ell_2(x)\pmod6
 \qquad\bigl(x\in(\mathbb Z/9\mathbb Z)^\times\bigr).
\]
La elección del menor generador positivo fija \(\ell_2=\ell_9\),
con los valores de \eqref{act21:tpk:logaritmo}.
\end{proposition}
\begin{proof}
Las potencias de \(2\) recorren \(1,2,4,8,7,5\) y cierran en el
sexto paso. Los generadores de un grupo cíclico de orden seis
corresponden a los exponentes coprimos con seis: \(1\) y \(5\).
Sus representantes son \(2\) y \(2^5\equiv5\pmod9\).
Como \(5\equiv2^{-1}\pmod9\), expresar una unidad como potencia de
\(5\) cambia el signo de su exponente en base \(2\). La
normalización por el menor representante positivo selecciona \(2\).
\end{proof}

La condición de homomorfismo por sí sola permite también aplicaciones
de imagen propia. La proposición distingue, por ello, la representación
fiel del grupo de unidades de sus cocientes. La extensión por cero
determina además el tratamiento de los tres residuos no invertibles;
el registro de procedencia conserva cuál de ellos fue observado.

La comparación de los cuatro observables recorre las mismas
\(104\,976\) condiciones iniciales y agrupa sus emisiones de seis
ventanas. Denotando por \(E_g\) el emisor con observable multiplicativo
\(g\), se obtiene el siguiente censo de imágenes.
\begin{table}[htbp]
\centering
\small
\begin{tabular}{@{}lrrrl@{}}
\toprule
Observable & \(|\operatorname{im}E_g|\) & Raíces nulas
 & Ley aditiva & Normalización\\
\midrule
\(\ell_2\)&468&sí&sí&menor generador\\
\(\ell_5\)&468&sí&sí&generador inverso\\
\(v_3\)&144&no&sí&imagen unitaria nula\\
\(\operatorname{id}\)&684&no&no&residuo positivo\\
\bottomrule
\end{tabular}
\caption{Comparación de cuatro calibres sobre un mismo transporte y un mismo dominio de condiciones iniciales.}
\label{mg05:tabla-calibres}
\end{table}

Los dos logaritmos satisfacen la ley aditiva de exponentes. La valuación
\(v_3\) vale cero en todas las unidades, pero toma los valores
\(1,1,2\) en \(3,6,9\); de ahí la diferencia entre sus dos columnas
de control. La identidad conserva el residuo positivo y falla la ley
del homomorfismo a \(\mathbb Z/6\mathbb Z\); por ejemplo, para
la unidad \(1\), la igualdad exigida sería \(1\equiv2\pmod6\).
El censo de \(468\) emisiones es común a ambos logaritmos y, por
tanto, conserva la ambigüedad de inversión. El criterio del menor
generador resuelve esa ambigüedad antes de cualquier lectura regional
de constantes. Los cardinales de la tabla describen la evaluación
exhaustiva de estos cuatro observables sobre el dominio fijado; su
función es contrastar la normalización operacional del emisor.
